\ifdefined\chinchin
	\documentclass[12pt, orivec]{extarticle} % 12pt, 14pt, 17pt, 20pt
\else
	\documentclass[10pt, orivec]{extarticle} % 12pt, 14pt, 17pt, 20pt
\fi
\usepackage{parskip}
\usepackage{amsmath}
\usepackage{amssymb}    % for \rightsquigarrow
\usepackage{wasysym}	% for frown face
\usepackage{mathrsfs} 	% for \mathscr
\usepackage{stmaryrd}
% \usepackage{unicode-math}
\usepackage[many]{tcolorbox}
\usepackage{ulem}
\usepackage{tikz-cd}		% commutative diagrams
\usepackage{tikz}
\usepackage{amsthm}
\usepackage{enumitem}	% for \ietemize custom labels
\usepackage{turnstile}	% longer turnstiles
\usepackage[sf,bf,big,raggedright,compact]{titlesec}	% change section color to blue
\usepackage{graphicx}	% Allows including images
% \usepackage[backend=biber,bibstyle=authoryear,citestyle=../authoryearbrack]{biblatex}
% \bibliography{../AGI-book}

\newtheorem{theorem}{Theorem}

\ifdefined\chinchin
	\newcommand{\cc}[2]{#1}
	\usepackage[CJKspace]{xeCJK}
	%\setCJKmainfont[BoldFont=SimHei,ItalicFont=AR PL KaitiM GB]{Alibaba PuHuiTi}
	\setCJKmainfont[Path={/usr/share/fonts/truetype/Alibaba PuHuiTi/}, BoldFont=Alibaba_PuHuiTi_Regular.ttf, ItalicFont=ukai.ttc]{Alibaba_PuHuiTi_Light.ttf}
	% \setmainfont[ItalicFont=Latin Modern Roman Slanted]{Alibaba Sans:style=Light}
\else
	\newcommand{\cc}[2]{#2}
	% \usepackage[no-math]{fontspec}
	% \setmainfont[Path={/usr/share/fonts/opentype/Alibaba Sans/},
	% BoldFont=Alibaba_Sans_Regular.otf,
	% ItalicFont=Alibaba_Sans_Light_Italic.otf]
	% {Alibaba_Sans_Light.otf}
	%\setmainfont{AlibabaSans_Light.otf}
	%\renewcommand{\baselinestretch}{0.8}
\fi

%\setlength{\headheight}{0cm}
%\setlength{\hoffset}{0cm}
%\setlength{\topmargin}{-2cm}
%\setlength{\oddsidemargin}{-2cm}
%\setlength{\evensidemargin}{-2cm}
%\setlength{\textwidth}{20cm}
%\setlength{\textheight}{26cm}
%\setlength{\headsep}{0.5em}
%\setlength{\topskip}{0.5em}
%\setlength{\footskip}{0.9cm}  % between bottom of page and page number
%\setlength{\floatsep}{0cm}
%\setlength{\textfloatsep}{0.6cm}
%\setlength{\intextsep}{0.5cm}
%\setlength{\parindent}{0em}   % em = width of capital M
%\setlength{\parskip}{10pt plus 5pt}

% Fix spilling of titles in bibliography:
%\DeclareFieldFormat*{title}{#1}
%
%\DeclareFieldFormat*{titlecase}{%
%    \ifdef{\currentfield}
%      {\ifcurrentfield{title}
%         {\usefield{\uline}{\currentfield}}%
%         {#1}}
%      {#1}}

\setcounter{secnumdepth}{1}		% 0 = no section numbers

\titleformat{\section}[hang]{\bfseries\Large\color{blue}}{\thesection \hspace{20pt}}{0pt}{}
\titleformat{\subsection}[hang]{\bfseries\large\color{blue}}{\thesubsection \hspace{10pt}}{0pt}{}
\titleformat{\subsubsection}[hang]{\bfseries\color{blue}}{}{0pt}{}

\itemsep1em
% \setlist[itemize]{noitemsep, topsep=15pt, partopsep=10pt}
\setlist[itemize]{topsep=15pt, partopsep=10pt}
\renewcommand{\labelitemi}{\textbullet}

\let\varzero\emptyset
\let\emptyset\varnothing		% round empty set
\newcommand{\vect}[1]{\symbf{#1}}
\newcommand{\book}[1]{$\NewSym[0.4]{../book-icon.png} \quad$ \parbox{0.9\textwidth}{\footnotesize #1}}
\newcommand{\code}[1]{{\small{\ttfamily #1}}}
\newcommand{\tab}{\hspace*{2cm}}
\newcommand{\powerset}{\raisebox{.15\baselineskip}{\Large\ensuremath{\wp}}}
\newcommand{\Chi}{\raisebox{2.5pt}{$\chi$}}
\newcommand*\KB{\vcenter{\hbox{\includegraphics{../KB-symbol.png}}}}
\newcommand*\NewSym[2][0.5]{\vcenter{\hbox{\includegraphics[scale=#1]{#2}}}}
\newcommand*\sigmoid{\vcenter{\hbox{\includegraphics{../sigmoid3.png}}}}
\newcommand{\smbox}[1]{\boxed{\footnotesize{\mbox{#1}}}}

\newcommand{\tikzmark}[1]{\tikz[overlay,remember picture] \node (#1) {};}

\newcommand{\Dfrac}[2]{%
\ooalign{%
      $\genfrac{}{}{2.9pt}0{\hphantom{#1}}{\hphantom{#2}}$\cr%
      $\color{white}\genfrac{}{}{1.5pt}0{\hphantom{#1}}{\hphantom{#2}}$\cr%
      $\color{white}\genfrac{}{}{1pt}0{\color{black}#1}{\color{black}#2}$}}

% \renewcommand{\thefootnote}{\fnsymbol{footnote}}
%\usepackage{perpage}
%\MakePerPage{footnote}
%\renewcommand{\thefootnote}{\ifcase\value{footnote}\or{*}
%	\or{$\dagger$}\or{**}\or{$\ddagger$}\fi}
%\renewcommand{\thefootnote}{\arabic{footnote}}
%\setcounter{footnote}{0}

\usepackage[hang,flushmargin]{footmisc}
\interfootnotelinepenalty=10000


\usepackage[left=1.5cm, right=1.5cm,top=1cm,bottom=1cm]{geometry}

\pagenumbering{gobble}

\begin{document}

\begin{itemize}

	\item[1] 永生是终极的幸福，对永生的追求 = 对至善的追求 \\
	因爲兩者的行爲似乎不谋而合 \\
	举例来说，《圣经 • 新约》约翰福音 3:16 \\
	被认为是基督教的中心思想： \\
	「神爱世人，甚至将祂的独生子献给世人， \\
	让信祂的人不至灭亡，反得永生」 \\
	我的意思并不是建议我们盲目信仰基督教 \\
	而是说，古人其实也思考永生与道德伦理的问题 \\
	他们的思想与现代思想往往高度重合 \\
	（但也不一定完全重合）

	\item[1] Eternal life is the ultimate happiness. \\
	Eternal life can be regarded as identical to the ``highest good'' \\
	in the sense that the pursuit of the two goals seem to coincide. \\
	For example, in the Christian Bible, John 3:16 is often considered \\
	the central idea of the New Testament or of Christianity: \\
	``For God so loved the world that he gave his only begotten Son, \\
	that whoever believes in him should not perish, but have eternal life.'' \\
	My point here is not to preach Christianity as scientific truth \\
	but rather to point out that the ancients have also pondered the questions \\
	of eternal life and morality \\
	and their ideas often overlap significantly with modern thinking \\
	(though they are not necessarily 100\% correct)

	\item[2] 纵观人类历史发展的规律， \\
	古文明总是被新兴的帝国征服、而趋向衰落 \\
	这似乎与热力学第二定律、熵的增加有关 \\
	但严格来说此定律只对封闭体系适用 \\
	而地球上的生物圈不是封闭体系（太阳给我们能量） \\
	主宰生命的大统一理论是 进化论 \\
	但它与热力学第二定律的关系，我们暂不清楚 \\
	（那牵涉所谓 “生命起源” 的问题） \\
	而，以进化论的观点来看， \\
	旧的、传统的思想与技术，被新的思想与技术取代 \\
	在这过程中，一个族裔 / 谱系 可以无止境地延续下去 \\
	在这意义上「衰老」和「死亡」并不是必然的 \\
	关键是要学习永生的秘诀，而这奥秘是不断变化的

	\item[2] Looking at the patterns of human history, \\
	older civilizations have always declined and been conquered by newer ones. \\
	This seems to be related to the second law of thermodynamics \\
	and the inexorable increase of entropy. \\
	However, strictly speaking, the 2nd law only applies to closed systems \\
	but the earth's biosphere is not a closed system (the sun gives us energy) \\
	The grand theory governing life is evolution \\
	but its relationship to the 2nd law is still unclear \\
	From an evolutionary perspective, \\
	older, traditional ideas and technologies are replaced by newer ones \\
	In this process, a people or lineage can continue its existence indefinitely \\
	Thus, ``decline'' and ``death'' are not inevitable \\
	The key is to master the secret of eternal life \\
	but this secret does not stay constant with time.

	\item[3] 再回到「帝国的兴起终将衰落」的轮回： \\
	从进化的观点来看，这并非一个必然的宿命 \\
	在殖民主义时期，欧洲人在非洲做了泯灭人性的破坏 \\
	时至今日，非洲仍然是祸乱不止的一片土地 \\
	新兴的霸权有时直接破坏较弱的群体、\\
	有时则阻止那些群体学习和进步 \\
	某意义上，这是符合进化论的「弱肉强食」 \\
	但另一角度看，群体或许厌倦了互相杀戮，而进化出合作和道德感 \\
	而我提出的理解就是： \\
	由此进化出来的道德，其轨迹就跟上面说的「至善」重合 \\
	我们只要找到「至善」的 pathway \\
	就可以达到永生，也就是所谓 “the survivor strain” \\
	即生物学上，群体经过很多代进化，仍然存活下来的那一条谱系

	\item[3] Returning to the seemingly inescapable cycle \\
	of the ``rise and fall of empires'' \\
	From an evolutionary perspective, this may not be an inevitable fate. \\
	Emerging empires sometimes directly attack and destroy weaker states \\
	while at other times prevent them from learning and progressing. \\
	During the colonial period, Europeans wreaked havoc in Africa. \\
	To this day, Africa remains a land of constant turmoil. \\
	In a sense, this kind of dog-eat-dog struggle is consistent with evolution. \\
	But from another perspective, \\
	the collective may tire of killing each other \\
	from which cooperation and morality could emerge. \\
	My understanding is that this evolutionary trajectory of morality \\
	would ultimately coincide with the ``highest good'' mentioned above. \\
	As long as we find the pathway to this ``highest good'' \\
	we can achieve immortality \\
	or in other words, we become the ``survivor strain'' \\
	ie, the lineage that survives each and every generation of evolution

	\item[4] 可以想象，随着时代进步， \\
	越往未来推进，则人类能获得生命延续的机会越大 \\
	那么，在现在这个看来仍然有点渺茫的时刻， \\
	人可以做什么去增加达到永生的可能呢？ \\
	踩低别人，尽量自私贪婪地赚钱？ \\
	我觉得还是应该尽快建立有兴趣追求永生的组织 \\
	这些组织既有平等（互相帮助）的观念， \\
	亦有保持进化的竞争力的观念 \\
	这跟人类社会很多现有的组织是不谋而合的： \\
	例如 教会、国家、家庭、企业 等 \\
	但我们要将目标更明确化、用更科学的方法制定群体的政策

	\item[4] It's conceivable that as time progresses \\
	the further into the future we advance \\
	the greater our chances of achieving life extension. \\
	So, at this moment, when the prospect of immortality still looks bleak \\
	what can we do to increase our chances? \\
	Do we trample on each other and act as selfishly as possible? \\
	I think we should organize ourselves towards the goal of immortality. \\
	Our organizations should value both egality \\
	as well as freedom (to maintain evolutionary competitiveness) \\
	This resonates with many existing organizations in human society \\
	such as churches, nations, families, and businesses \\
	We just need to make clearer our purpose \\
	and to design our organizational policies more scientifically

\end{itemize}

\end{document}
